\documentclass{article}
\usepackage[utf8]{inputenc}
\usepackage{graphicx}
\usepackage{hyperref}
\usepackage{listings}
\usepackage{geometry}
\geometry{a4paper, margin=1in}

\title{Preprocessing Data Volume Pengangkutan Sampah}
\author{Rayhan Surya Destian}
\date{\today}

\begin{document}

\maketitle

\section{Langkah-Langkah Pra-Pemrosesan Data}

\subsection{Memuat Dataset}
Dataset dibaca ke dalam DataFrame menggunakan Pandas dengan delimiter titik koma (\texttt{;}).

\subsection{Memeriksa Data}
Dilakukan inspeksi awal menggunakan \texttt{df.head()}, \texttt{df.info()}, dan \texttt{df.isnull().sum()} untuk memahami struktur data dan mengidentifikasi nilai yang hilang.

\subsection{Menangani Inkonsistensi Satuan}
Kolom \texttt{satuan\_panjang/luas} diperiksa untuk memastikan konsistensi. Nilai yang hilang diisi dengan 'm' (meter) berdasarkan konteks.

\subsection{Membersihkan Kolom Volume Sampah}
Entri tidak valid seperti '-' diganti dengan \texttt{NaN} dan dikonversi menjadi tipe data numerik.

\subsection{Mengisi Nilai yang Hilang}
\begin{itemize}
    \item \textbf{Volume Sampah}: Nilai yang hilang diisi dengan rata-rata volume per \texttt{kecamatan}.
    \item \textbf{Panjang/Luas}: Nilai yang hilang diisi dua kali; pertama dengan rata-rata per \texttt{kecamatan}, kemudian dengan rata-rata keseluruhan jika masih ada yang kosong.
\end{itemize}

\subsection{Mengonversi Tipe Data}
Pastikan semua kolom memiliki tipe data yang sesuai, seperti integer untuk \texttt{bulan} dan \texttt{tanggal}, serta float untuk \texttt{panjang/luas} dan \texttt{volume\_sampah\_perhari(m3)}.

\subsection{Pemeriksaan Akhir}
Verifikasi bahwa tidak ada nilai yang hilang tersisa menggunakan \texttt{df.isnull().sum()}.

\section{Diagram Alur Pra-Pemrosesan Data}

\begin{figure}[h!]
\centering
\includegraphics[width=0.3\linewidth]{preprocessing_diagram.png}
\caption{Alur Pra-Pemrosesan Data}
\end{figure}

\end{document}
